\subsection{Simulated annealing}
\label{subsec:simulated-annealing}

Gradient descent and its stochastic variants are fundamentally local search methods.
They can move quickly once they are inside a good basin, but on objectives with several separated local minima, saddles, or energy barriers they do not provide a reliable global-optimization strategy.
One response is to use careful initialization; another is brute force, i.e.\ many dispersed restarts followed by local refinement.
Simulated annealing takes a different approach:
it combines the accept--reject mechanism of Metropolis--Hastings with a temperature schedule that is high at the beginning and small later on.
Here the auxiliary mechanism is the scalar temperature itself.
At iteration $k$ the method works on an enlarged description $(x,T_k)$ rather than on the state $x$ alone: the temperature temporarily changes the landscape, then is cooled back toward the original optimization problem.
The high-temperature phase makes it easier to leave poor local minima, while the low-temperature phase gradually concentrates the search near low-energy states.

\subsubsection{Boltzmann targets and temperature flattening}
\label{subsubsec:sa-tempered-targets}

Let $f:\mathcal{X}\to\R$ be the objective we want to minimize.
For each temperature $T>0$, define the Gibbs or Boltzmann distribution
\begin{equation}
  \pi_T(x)
  \coloneqq
  \frac{\exp(-f(x)/T)}{Z_T},
  \qquad
  Z_T \coloneqq \int_{\mathcal{X}} \exp(-f(x)/T)\,dx,
  \label{eq:sa-boltzmann}
\end{equation}
with the integral replaced by a sum when $\mathcal{X}$ is discrete.
The theoretical probabilistic interpretation requires $Z_T<\infty$, but the algorithm itself only uses ratios of the form $\pi_T(y)/\pi_T(x)$, so the normalizing constant never needs to be computed explicitly.

Two simple facts explain why \eqref{eq:sa-boltzmann} is useful.
First, for every $T>0$, maximizing $\pi_T(x)$ is exactly the same as minimizing $f(x)$ because the map $u\mapsto e^{-u/T}$ is strictly decreasing.
Second, temperature controls how strongly energy differences are emphasized.
When $T$ is large, the factor $e^{-f(x)/T}$ changes slowly with $f(x)$, so the target is relatively flat.
When $T$ is small, differences in objective value are amplified and the mass concentrates around the lowest-energy states.

\begin{figure}[t]
  \centering
  \begin{tikzpicture}
  \begin{axis}[
    width=0.9\linewidth,
    height=6.4cm,
    xmin=-2.6, xmax=2.8,
    ymin=0, ymax=1.05,
    axis lines=left,
    xlabel={$x$},
    ylabel={relative density},
    xtick=\empty,
    ytick={0,0.5,1.0},
    yticklabels={$0$,$0.5$,$1$},
    clip=false,
    legend style={
      draw=stat221Gray!25,
      fill=white,
      fill opacity=0.95,
      text opacity=1,
      font=\small,
      cells={anchor=west},
      at={(0.98,0.98)},
      anchor=north east,
    },
  ]
    \addplot[stat221Gray, densely dashed, thick, samples=220, domain=-2.6:2.8]
      {0.10 + 0.18*(0.18*(x^2-4)^2 + 0.25*x + 0.55)};
    \addlegendentry{rescaled objective}

    \addplot[stat221Red, thick, samples=220, domain=-2.6:2.8]
      {exp(-((0.18*(x^2-4)^2 + 0.25*x + 0.505))/1.6)};
    \addlegendentry{hot: $T=1.6$}

    \addplot[stat221Blue, thick, samples=220, domain=-2.6:2.8]
      {exp(-((0.18*(x^2-4)^2 + 0.25*x + 0.505))/0.7)};
    \addlegendentry{warm: $T=0.7$}

    \addplot[stat221Purple, thick, samples=220, domain=-2.6:2.8]
      {exp(-((0.18*(x^2-4)^2 + 0.25*x + 0.505))/0.25)};
    \addlegendentry{cold: $T=0.25$}

    \node[font=\scriptsize\color{stat221Gray}, anchor=north] at (axis cs:-2.03,0.09) {global well};
    \node[font=\scriptsize\color{stat221Gray}, anchor=south] at (axis cs:1.95,0.26) {local well};
  \end{axis}
\end{tikzpicture}

  \caption{For a fixed two-well objective (shown faintly in gray after rescaling), the quantities proportional to $\exp(-f(x)/T)$ are much flatter when $T$ is large and much more concentrated near low-energy states when $T$ is small.
  Simulated annealing exploits this by exploring broadly while the system is hot and becoming increasingly selective as it cools.}
  \label{fig:sa-temperature-flattening}
\end{figure}

Thinking of $T$ as a physical temperature gives the right intuition:
when the system is hot, the wells of the landscape are relatively shallow, so movement between them is easier;
when the system is cold, those same wells become deep, and the dynamics become much more reluctant to move uphill.
\Cref{fig:sa-temperature-flattening} makes this concrete: the hot distribution still assigns substantial mass to the higher-energy local well and even to the barrier region, whereas the cold distribution is sharply concentrated near the global well.

\subsubsection{Metropolis updates and simulated annealing}
\label{subsubsec:sa-mh-view}

Suppose that at a fixed temperature $T$ we use a proposal kernel $q(y\mid x)$.
If we were running an ordinary Metropolis--Hastings chain targeting $\pi_T$, then the acceptance probability would be
\begin{equation}
  \alpha_T(x,y)
  =
  \min\left\{
    1,\,
    \exp\!\left(-\frac{f(y)-f(x)}{T}\right)\frac{q(x\mid y)}{q(y\mid x)}
  \right\}.
  \label{eq:sa-mh-ratio}
\end{equation}
For a symmetric random-walk proposal, $q(y\mid x)=q(x\mid y)$, so \eqref{eq:sa-mh-ratio} simplifies to
\begin{equation}
  \alpha_T(x,y)
  =
  \min\left\{
    1,\,
    \exp\!\left(-\frac{f(y)-f(x)}{T}\right)
  \right\}.
  \label{eq:sa-symmetric-acceptance}
\end{equation}
This form is especially transparent:
downhill moves ($f(y)\le f(x)$) are always accepted, while uphill moves are accepted with a probability that decays exponentially in the energy increase $f(y)-f(x)$.
Large temperature makes uphill moves fairly likely; small temperature suppresses them.
Temperature turns uphill acceptance into a controllable exploration mechanism: hot chains are willing to climb barriers, while cold chains behave much more like descent.

Simulated annealing uses the same accept--reject idea, but with a nonincreasing sequence of temperatures
\[
  T_0 \ge T_1 \ge T_2 \ge \cdots \downarrow 0.
\]
At iteration $k$, the algorithm updates the pair $(x^{(k)},T_k)$ by holding the temperature fixed at its current value and performing a Metropolis--Hastings step for the target $\pi_{T_k}$.
Because the temperature changes over time, the chain is no longer trying to sample from one fixed distribution.
Instead, it gradually transitions from exploratory dynamics at high temperature to almost-greedy local search at low temperature.

\begin{algorithm}
  \caption{Random-walk simulated annealing}
  \label{alg:simulated-annealing}
  \begin{algorithmic}[1]
    \Require Objective $f$, symmetric proposal kernel $q(\cdot\mid x)$, temperatures $T_0,\dots,T_{K-1}$ with $T_{k+1}\le T_k$, initial state $x^{(0)}$
    \Ensure Best state encountered during the run
    \State $x_{\mathrm{best}} \gets x^{(0)}$, $f_{\mathrm{best}} \gets f(x^{(0)})$
    \For{$k=0,1,\dots,K-1$}
      \State Sample a proposal $y \sim q(\cdot\mid x^{(k)})$
      \State $\alpha_k \gets \min\!\left\{1,\exp\!\left(-(f(y)-f(x^{(k)}))/T_k\right)\right\}$
      \State Draw $U_k \sim \mathrm{Unif}[0,1]$
      \If{$U_k \le \alpha_k$}
        \State $x^{(k+1)} \gets y$
      \Else
        \State $x^{(k+1)} \gets x^{(k)}$
      \EndIf
      \If{$f(x^{(k+1)}) < f_{\mathrm{best}}$}
        \State $x_{\mathrm{best}} \gets x^{(k+1)}$, $f_{\mathrm{best}} \gets f(x^{(k+1)})$
      \EndIf
    \EndFor
    \State \Return $x_{\mathrm{best}}$
  \end{algorithmic}
\end{algorithm}

The best-so-far bookkeeping is important.
Even when the current state wanders away from a good point, the algorithm retains the lowest objective value it has seen so far.
For asymmetric proposals one simply replaces the acceptance formula in \cref{alg:simulated-annealing} by the full Metropolis--Hastings ratio \eqref{eq:sa-mh-ratio}.
As $T_k$ becomes small, the chain behaves more and more like randomized local descent on $f$, but the early hot phase is what gives the method a chance to reach a good basin before that local behavior takes over.

\subsubsection{Cooling schedules and global guarantees}
\label{subsubsec:sa-schedules}

If one wants theorem-level global guarantees rather than a practical heuristic, the cooling schedule must be extremely slow.
A classical theorem-level schedule is logarithmic cooling of the form
\begin{equation}
  T_k = \frac{c}{\log(k+k_0)},
  \qquad c>0,
  \label{eq:sa-logarithmic}
\end{equation}
in the finite-state Gibbs setting analyzed by \citet[eq.~(12.7), Thm.~C, and p.~732]{geman_geman_1984_stochastic_relaxation}.
In that theorem, the constant $c$ is problem-dependent and reflects the barrier structure of the energy landscape.
These schedules are slow precisely because they are designed to give the algorithm enough time to escape increasingly deep local wells.
We are using \eqref{eq:sa-logarithmic} only as a theorem-level schematic, not as a drop-in prescription: outside that finite-state setting, the exact hypotheses and constants depend delicately on the proposal kernel and the geometry of the state space.
Its message is structural: provable global exploration is much slower than the schedules people usually want to run.

\subsubsection{Practical schedules and proposal design}
\label{subsubsec:sa-practical}

The proposal distribution matters just as much here as it does in ordinary Metropolis--Hastings.
For continuous state spaces, Gaussian random-walk proposals are a common default.
On structured discrete spaces, one instead designs local combinatorial moves that preserve validity of the state, such as swaps, reversals, or Gibbs-style coordinate updates.

In practical computation one often uses faster heuristics such as a linear schedule
\begin{equation}
  T_k = T_{\max} - \frac{k}{K}\bigl(T_{\max}-T_{\min}\bigr),
  \label{eq:sa-linear}
\end{equation}
or an exponential schedule
\begin{equation}
  T_k = T_{\max}\alpha^k,
  \qquad \alpha \in (0,1),
  \label{eq:sa-exponential}
\end{equation}
because they cool more aggressively.
These schedules are heuristics: they often work far better computationally than logarithmic cooling, but they do not come with the same global guarantee.
Cool too quickly and the method behaves like a greedy local search that may freeze in a poor basin; cool more slowly and the method explores better, but at higher computational cost.
The temperature schedule is therefore the main dial that trades global exploration against local refinement.
Because the method is stochastic and still sensitive to initialization, one often runs several independent seeds and compares the best objective values they find rather than treating a single annealing trajectory as definitive.
That comparison separates two questions that are easy to conflate in a single run: whether one schedule explores a basin effectively once it reaches it, and whether the overall procedure reliably reaches good basins at all.

This tension between hot exploration and cold exploitation is the bridge to tempering in the next unit.
Simulated annealing commits to a single cooling schedule and eventually gives up the hot dynamics.
Tempering keeps a finite ladder of hot auxiliary distributions alive throughout the run and couples them to the cold target, so barrier-crossing remains available even after the chain has returned to the original distribution; see \Cref{subsec:tempering}.

\begin{example}[Traveling salesman and a 2-opt proposal]
\label{ex:sa-tsp}
Given locations $z_1,\dots,z_n\in\R^2$, a tour is a permutation $\sigma=(i_1,\dots,i_n)$ of $\{1,\dots,n\}$.
The traveling-salesman objective is
\[
  f(\sigma)
  =
  \sum_{j=1}^{n} \norm{z_{i_j}-z_{i_{j+1}}}_2,
  \qquad i_{n+1}\coloneqq i_1,
\]
and the goal is to minimize $f(\sigma)$ over all tours.
This problem is NP-hard, so exhaustive search is generally impossible.

A standard symmetric proposal is the \emph{2-opt} move:
choose two indices $a<b$ and reverse the segment $(i_a,i_{a+1},\dots,i_b)$.
This keeps the tour valid while making a large enough local change to repair bad crossings or reorganize an awkward portion of the route.
Because the move is symmetric, the acceptance probability is exactly \eqref{eq:sa-symmetric-acceptance}.
Early in the run the algorithm is willing to accept a longer tour if that helps reorganize the route; later, once the temperature is small, it behaves much more like a greedy local improvement method.

The handwritten comparison on one fixed TSP instance made the schedule tradeoff concrete.
A very short aggressive schedule froze quickly and often kept visibly poor crossings, because the chain stopped accepting uphill reorganizations before the 2-opt proposals had time to rearrange the route.
A longer linear schedule improved the tour more reliably, and a slower exponential schedule typically produced the cleanest final route on that example.
The point was not that one named schedule always dominates.
It was that schedule length has to be matched to the proposal's ability to make useful large-scale repairs before the search becomes effectively greedy.
\end{example}
